\documentclass[11pt,a4paper]{article}
\usepackage[T1]{fontenc}
\usepackage[utf8]{inputenc}
\usepackage[margin=1in]{geometry}
\usepackage{amsmath}
\usepackage{amssymb}
\usepackage{hyperref}

\title{Zero-Effort Miss (ZEM) and Its Relation to Proportional Navigation}
\author{PX4 Python SITL Guidance Notes}
\date{\today}

\begin{document}
\maketitle

\section{Variable Inventory}
\begin{center}
\renewcommand{\arraystretch}{1.2}
\begin{tabular}{|c|p{7.1cm}|c|c|}
\hline
\textbf{Symbol} & \textbf{Description / Definition} & \textbf{Unit} & \textbf{Dimension} \\
\hline
$\mathbf{r}_M$ & Interceptor position vector & m & $L$ \\
\hline
$\mathbf{r}_T$ & Target position vector & m & $L$ \\
\hline
$\mathbf{v}_M$ & Interceptor velocity vector & m/s & $LT^{-1}$ \\
\hline
$\mathbf{v}_T$ & Target velocity vector & m/s & $LT^{-1}$ \\
\hline
$\mathbf{r}$ & Relative position, $\mathbf{r}=\mathbf{r}_T-\mathbf{r}_M$ & m & $L$ \\
\hline
$\mathbf{v}$ & Relative velocity, $\mathbf{v}=\mathbf{v}_T-\mathbf{v}_M$ & m/s & $LT^{-1}$ \\
\hline
$R$ & Range, $R=\lVert \mathbf{r} \rVert$ & m & $L$ \\
\hline
$\hat{\mathbf{r}}$ & Line-of-sight unit vector, $\hat{\mathbf{r}}=\mathbf{r}/\lVert\mathbf{r}\rVert$ & 1 & 1 \\
\hline
$V_c$ & Closing speed, $V_c=-\hat{\mathbf{r}}^\top\mathbf{v}$ (positive when target and interceptor are closing) & m/s & $LT^{-1}$ \\
\hline
$t_{go}$ & Time-to-go approximation, $t_{go}\approx R/V_c$ & s & $T$ \\
\hline
$\mathbf{ZEM}$ & Zero-Effort Miss vector, $\mathbf{ZEM}=\mathbf{r}+\mathbf{v}t_{go}$ & m & $L$ \\
\hline
$\mathbf{ZEM}_{\perp}$ & Lateral ZEM, $\mathbf{ZEM}_{\perp}=(\mathbf{I}-\hat{\mathbf{r}}\hat{\mathbf{r}}^\top)\mathbf{ZEM}$ & m & $L$ \\
\hline
$\mathbf{a}_c$ & Commanded acceleration from ZEM guidance, $\mathbf{a}_c=N\mathbf{ZEM}_{\perp}/t_{go}^2$ & m/s$^2$ & $LT^{-2}$ \\
\hline
$N$ & Navigation gain & 1 & 1 \\
\hline
$\lambda$ & LOS angle (planar PN form) & rad & 1 \\
\hline
$\dot{\lambda}$ & LOS angular rate & rad/s & $T^{-1}$ \\
\hline
$a_n$ & PN normal acceleration, $a_n=NV_c\dot{\lambda}$ & m/s$^2$ & $LT^{-2}$ \\
\hline
\end{tabular}
\end{center}

\section{Coordinate Systems and Command Frame}
This note uses a navigation/inertial frame for engagement geometry and a body frame for vehicle actuation:
\begin{itemize}
    \item Navigation frame $\{n\}$: NED (North-East-Down), where positions and relative vectors are formed.
    \item Body frame $\{b\}$: FRD (Forward-Right-Down), where aerodynamic and actuator commands are typically applied.
\end{itemize}

Let $R_{bn}$ be the direction cosine matrix from navigation to body, and $R_{nb}=R_{bn}^\top$ from body to navigation.
For any vector $\mathbf{x}$,
\begin{equation}
\mathbf{x}_b = R_{bn}\mathbf{x}_n, \qquad \mathbf{x}_n = R_{nb}\mathbf{x}_b.
\end{equation}

\section{Setup and Notation}
\subsection{3D Vector Form}
Let the interceptor be denoted by subscript $M$ and the target by subscript $T$.
Define relative position and velocity as
\begin{align}
\mathbf{r} &= \mathbf{r}_T - \mathbf{r}_M, \\
\mathbf{v} &= \mathbf{v}_T - \mathbf{v}_M,
\end{align}
with range and LOS unit vector
\begin{align}
R &= \lVert \mathbf{r} \rVert, \\
\hat{\mathbf{r}} &= \frac{\mathbf{r}}{\lVert \mathbf{r} \rVert}.
\end{align}

An often-used time-to-go approximation is
\begin{align}
V_c &= -\hat{\mathbf{r}}^\top \mathbf{v}, \\
t_{go} &\approx \frac{R}{V_c}, \qquad V_c > 0.
\end{align}

\subsection{2D Planar Form}
In planar engagement, use
\begin{equation}
\mathbf{r} = \begin{bmatrix} x \\ y \end{bmatrix}, \qquad
\mathbf{v} = \begin{bmatrix} \dot{x} \\ \dot{y} \end{bmatrix}, \qquad
R = \sqrt{x^2+y^2}, \qquad
\hat{\mathbf{r}} = \frac{1}{R}\begin{bmatrix} x \\ y \end{bmatrix}.
\end{equation}
The same formulas for $V_c$ and $t_{go}$ apply.

\section{Zero-Effort Miss (ZEM)}
\subsection{3D Vector Form}
The Zero-Effort Miss is the predicted miss vector at intercept time if no additional control is applied:
\begin{equation}
\mathbf{ZEM} = \mathbf{r} + \mathbf{v}\,t_{go}.
\end{equation}

Guidance uses the lateral component of ZEM, i.e. the component perpendicular to LOS:
\begin{equation}
\mathbf{ZEM}_{\perp} = \left(\mathbf{I} - \hat{\mathbf{r}}\hat{\mathbf{r}}^\top\right)\mathbf{ZEM}.
\end{equation}

The matrix
\begin{equation}
\mathbf{P}_{\perp} = \mathbf{I} - \hat{\mathbf{r}}\hat{\mathbf{r}}^\top
\end{equation}
is an orthogonal projector onto the plane normal to LOS. To see this, first note that
\begin{equation}
\mathbf{P}_{\parallel} = \hat{\mathbf{r}}\hat{\mathbf{r}}^\top
\end{equation}
projects any vector onto LOS, since
\begin{equation}
\mathbf{P}_{\parallel}\mathbf{x} = (\hat{\mathbf{r}}^\top\mathbf{x})\hat{\mathbf{r}}.
\end{equation}
Therefore,
\begin{equation}
\mathbf{P}_{\perp}\mathbf{x} = \mathbf{x} - (\hat{\mathbf{r}}^\top\mathbf{x})\hat{\mathbf{r}}
\end{equation}
removes the LOS-parallel component and leaves only the lateral component.

Applied to ZEM, this gives the decomposition
\begin{equation}
\mathbf{ZEM} = \underbrace{(\hat{\mathbf{r}}^\top\mathbf{ZEM})\hat{\mathbf{r}}}_{\text{parallel to LOS}} + \underbrace{\mathbf{ZEM}_{\perp}}_{\text{lateral to LOS}}.
\end{equation}
Also, by construction,
\begin{equation}
\hat{\mathbf{r}}^\top\mathbf{ZEM}_{\perp} = 0,
\end{equation}
so $\mathbf{ZEM}_{\perp}$ is exactly orthogonal to LOS.

Guidance uses $\mathbf{ZEM}_{\perp}$ because this is the part that creates cross-range miss. Steering acceleration is used to reduce this lateral miss to zero; the along-LOS part is mainly associated with closing/timing rather than geometric miss angle.

A common ZEM command is
\begin{equation}
\mathbf{a}_c = N\,\frac{\mathbf{ZEM}_{\perp}}{t_{go}^2},
\end{equation}
where $N$ is a dimensionless gain (navigation ratio).

\subsection{2D Planar Form}
In 2D, the same equations hold with 2-vectors. If desired, define the LOS-normal unit vector as
\begin{equation}
\hat{\mathbf{n}} = \begin{bmatrix} -\hat{r}_y \\ \hat{r}_x \end{bmatrix},
\end{equation}
then the scalar lateral miss is $z_{\perp}=\hat{\mathbf{n}}^\top\mathbf{ZEM}$ and the lateral acceleration command magnitude is approximately
\begin{equation}
a_{\perp} = N\,\frac{z_{\perp}}{t_{go}^2}.
\end{equation}

\section{Relation to Proportional Navigation (PN)}
\subsection{2D Planar PN}
Classical PN in planar form commands normal acceleration
\begin{equation}
a_n = N V_c \dot{\lambda},
\end{equation}
where $\dot{\lambda}$ is LOS angular rate.

Under standard assumptions (small LOS angle, nearly constant closing speed),
\begin{equation}
\lVert \mathbf{ZEM}_{\perp} \rVert \approx R\,\dot{\lambda}\,t_{go}.
\end{equation}
Substitute this into ZEM guidance:
\begin{align}
\lVert \mathbf{a}_c \rVert
&\approx N\frac{R\dot{\lambda}t_{go}}{t_{go}^2}
= N\frac{R}{t_{go}}\dot{\lambda}
\approx N V_c \dot{\lambda}.
\end{align}
Therefore, ZEM guidance reduces to PN in this regime.

\subsection{3D PN Interpretation}
In 3D, PN is often written with LOS angular-rate vector
\begin{equation}
\boldsymbol{\omega}_{LOS} = \frac{\mathbf{r} \times \mathbf{v}}{R^2},
\end{equation}
and acceleration command direction
\begin{equation}
\mathbf{a}_{PN} = N V_c\,(\boldsymbol{\omega}_{LOS} \times \hat{\mathbf{r}}).
\end{equation}
This command is perpendicular to LOS and is the direct 3D analogue of planar $a_n = N V_c\dot{\lambda}$.

\section{Stationary Target Special Case}
For a stationary target, $\mathbf{v}_T=\mathbf{0}$ and therefore
\begin{equation}
\mathbf{v} = -\mathbf{v}_M.
\end{equation}
The same ZEM equations apply directly.

In this repository, the script
\texttt{src/guidance/zero\_effort\_miss\_demo.py}
implements exactly this case in the North-East plane,
computes $\mathbf{ZEM}$ each step, extracts $\mathbf{ZEM}_{\perp}$,
and animates the interceptor trajectory against a stationary target.

\section{Practical Considerations}
In ZEM guidance as written here,
\begin{equation}
\mathbf{a}_{c,n} = N\,\frac{\mathbf{ZEM}_{\perp,n}}{t_{go}^2},
\end{equation}
so the commanded acceleration is naturally formed in the navigation/engagement frame.

\subsection*{Practical time-to-go estimate}
For ZEM, a practical time-to-go estimate is
\begin{equation}
V_c = -\hat{\mathbf{r}}^\top\mathbf{v},
\end{equation}
\begin{equation}
t_{go} = \mathrm{clamp}\!\left(\frac{R}{\max(V_c,\epsilon)},\ t_{go,min},\ t_{go,max}\right),
\end{equation}
with small $\epsilon > 0$ and explicit bounds to avoid singular or overly aggressive commands.
For a stationary target ($\mathbf{v}_T=\mathbf{0}$), this reduces to
\begin{equation}
\mathbf{v} = -\mathbf{v}_M,
\qquad
V_c = -\hat{\mathbf{r}}^\top(-\mathbf{v}_M) = \hat{\mathbf{r}}^\top\mathbf{v}_M.
\end{equation}

\subsection*{Can $\mathbf{a}_c$ be used directly as a body acceleration command?}
Usually \textbf{no, not directly}. You should first ensure both vectors are expressed in the same frame.
If your inner loop expects body-frame acceleration commands, convert first:
\begin{equation}
\mathbf{a}_{c,b} = R_{bn}\,\mathbf{a}_{c,n}.
\end{equation}

In practice, many autopilots do not track full body acceleration directly; they track attitude/rates and thrust.
So $\mathbf{a}_{c,n}$ (or $\mathbf{a}_{c,b}$) is commonly mapped into achievable attitude and thrust commands with saturation and vehicle-dynamics constraints.

\subsection{Example: PX4 Fixed-Wing (Bank-to-Turn) Realization}
Assume a PX4-based fixed-wing vehicle using coordinated bank-to-turn behavior, with available attitude and/or body-rate setpoints.

\paragraph{Step 1: Guidance law in navigation frame}
Compute either ZEM or PN lateral acceleration command in navigation frame:
\begin{align}
\mathbf{a}_{cmd,n}^{ZEM} &= N\frac{\mathbf{ZEM}_{\perp,n}}{t_{go}^2}, \\
\mathbf{a}_{cmd,n}^{PN}  &= N V_c\left(\boldsymbol{\omega}_{LOS}\times\hat{\mathbf{r}}\right),
\quad
\boldsymbol{\omega}_{LOS}=\frac{\mathbf{r}\times\mathbf{v}}{\|\mathbf{r}\|^2}.
\end{align}
Use only the lateral (cross-track) component for bank-to-turn.

\paragraph{Step 2: Convert lateral acceleration to roll command}
For coordinated turns (small sideslip), lateral acceleration is approximated by
\begin{equation}
a_{lat} \approx g \tan\phi,
\end{equation}
thus
\begin{equation}
\phi_{cmd} = \mathrm{sat}\!\left(\arctan\!\frac{a_{lat,cmd}}{g},\;-\phi_{max},\;+\phi_{max}\right).
\end{equation}
Equivalent turn-rate form:
\begin{equation}
\dot{\chi}_{cmd} \approx \frac{g}{V}\tan\phi_{cmd}
\quad\Longleftrightarrow\quad
\phi_{cmd}\approx\arctan\!\left(\frac{V\dot{\chi}_{cmd}}{g}\right).
\end{equation}

\paragraph{Step 3: Send PX4-compatible setpoints}
Two common options:
\begin{itemize}
    \item Attitude-setpoint path: command $\phi_{cmd}$ (roll), while pitch/throttle come from speed/altitude loops.
    \item Rate-setpoint path: inner roll-rate loop tracks $p_{cmd}$ from a roll-angle outer loop, e.g. $p_{cmd}=k_{\phi}(\phi_{cmd}-\phi)$ with rate saturation.
\end{itemize}

\paragraph{Vertical channel (pitch/throttle) for fixed-wing}
Bank-to-turn guidance primarily controls lateral geometry. Vertical motion is usually handled by a separate longitudinal loop:
\begin{itemize}
    \item use altitude / flight-path-angle / vertical-speed control to generate $\theta_{cmd}$;
    \item use throttle (energy control) to maintain airspeed and climb/descent capability.
\end{itemize}
In other words, lateral guidance gives roll demand, while longitudinal control gives pitch and thrust demand. If full 3D guidance acceleration is available, split it into lateral and vertical components, then map vertical demand into feasible $\theta_{cmd}$ and throttle with limits.

\paragraph{PI example for generating $q_{cmd}$ from $a_{z,cmd}$}
If you want to command body vertical acceleration directly (FRD body axis), a practical outer-loop PI law is
\begin{equation}
e_{az} = a_{z,cmd} - a_{z,meas},
\end{equation}
\begin{equation}
q_{cmd} = \mathrm{sat}\!\left(k_{p,az}e_{az} + k_{i,az}\eta + q_{ff},\ -q_{max},\ +q_{max}\right),
\end{equation}
with integrator state
\begin{equation}
\dot{\eta} = e_{az}.
\end{equation}

Recommended implementation details:
\begin{itemize}
    \item low-pass filter $a_{z,meas}$ before feedback;
    \item anti-windup: only integrate when not saturating in the same error direction;
    \item clamp integrator state $\eta$ to finite bounds;
    \item gate/reduce integrator near stall or very low airspeed.
\end{itemize}

One simple anti-windup rule is: if $q_{cmd}$ saturates high and $e_{az}>0$, pause integration; if $q_{cmd}$ saturates low and $e_{az}<0$, pause integration; otherwise integrate normally.

\paragraph{Direct roll-rate command from body-frame acceleration direction}
If commanded acceleration is already available in body FRD as
\begin{equation}
\mathbf{a}_{cmd,b} = \begin{bmatrix} a_x & a_y & a_z \end{bmatrix}^\top,
\end{equation}
and the design goal is to drive lateral component $a_y \rightarrow 0$ by roll-rate command directly, one practical heuristic is
\begin{equation}
p_{cmd} = \mathrm{sat}\!\left(-k_p\,\mathrm{atan2}(a_y, a_z),\ -p_{max},\ +p_{max}\right),
\end{equation}
where $k_p$ has units $\mathrm{s}^{-1}$ and the sign should be verified against the adopted FRD and positive-roll convention.

This avoids explicit roll-angle command and directly uses acceleration-vector tilt in the $y$-$z$ plane. In implementation, use filtering and deadband to reduce chatter near $a_y\approx 0$.

\paragraph{Step 4: Practical limits and scheduling}
Apply:
\begin{itemize}
    \item acceleration, bank-angle, and roll-rate saturation;
    \item low-pass filtering of LOS rate / ZEM terms;
    \item gain scheduling with airspeed $V$ (same $a_{lat}$ implies different turn-rate authority at different $V$);
    \item minimum $t_{go}$ floor for ZEM to avoid aggressive commands near intercept.
\end{itemize}

\paragraph{Interpretation}
For fixed-wing bank-to-turn, guidance mainly shapes lateral acceleration demand, then autopilot loops realize it through bank angle (and indirectly yaw/heading change), while thrust and pitch loops maintain energy and flight-path constraints.

\paragraph{Cycle-level pseudocode}
\begin{verbatim}
Given target state (r_T, v_T), ownship state (r_M, v_M), attitude R_bn:

1) Relative geometry:
   r = r_T - r_M
   v = v_T - v_M
   R = ||r||
   rhat = r / max(R, eps)
   Vc = -dot(rhat, v)

2) Guidance command (choose one):
   ZEM: tgo = clamp(R / max(Vc, eps), tgo_min, tgo_max)
        ZEM = r + v * tgo
        ZEM_perp = (I - rhat*rhat^T) * ZEM
        a_cmd_n = N * ZEM_perp / (tgo*tgo)
   PN:  omega_LOS = cross(r, v) / max(R*R, eps)
        a_cmd_n = N * Vc * cross(omega_LOS, rhat)

3) Extract lateral acceleration demand:
   a_lat_cmd = project_to_horizontal_lateral(a_cmd_n)

4) Convert to roll command:
   phi_cmd = sat(atan(a_lat_cmd / g), -phi_max, +phi_max)

5) Command PX4:
   attitude mode -> send roll setpoint phi_cmd
   or rate mode  -> p_cmd = sat(k_phi*(phi_cmd - phi), -p_max, +p_max)

6) Vertical channel (separate longitudinal loop):
   e_az = az_cmd - az_meas_lpf
   q_unsat = kp_az*e_az + ki_az*eta + q_ff
   q_cmd = sat(q_unsat, -q_max, +q_max)
   if not saturating_in_same_error_direction:
       eta = clamp(eta + e_az*dt, -eta_max, +eta_max)
   throttle_cmd = energy_or_speed_controller(V_cmd, V, h_cmd, h)

7) Optional direct-rate heuristic from body acceleration:
   p_cmd = sat(-k_p * atan2(a_cmd_b.y, a_cmd_b.z), -p_max, +p_max)

8) Apply filters, saturation, and airspeed scheduling each cycle.
\end{verbatim}

\section{Camera-Based Guidance with Forward-Looking FRD X-Axis}
Assume an onboard camera with optical axis aligned to body forward axis in FRD:
\begin{equation}
\hat{z}_{optical} \parallel +X_b.
\end{equation}
Define camera frame $\{c\}$ and body frame $\{b\}$. If camera is perfectly aligned with body,
\begin{equation}
R_{bc}=I,
\end{equation}
otherwise use calibrated extrinsics $R_{bc}$.

\subsection{From image detection to LOS unit vector}
For a detected target pixel $(u,v)$ with intrinsics matrix $K$, form
\begin{equation}
\tilde{\boldsymbol{\ell}}_c = K^{-1}\begin{bmatrix}u & v & 1\end{bmatrix}^\top,
\qquad
\hat{\boldsymbol{\ell}}_c = \frac{\tilde{\boldsymbol{\ell}}_c}{\|\tilde{\boldsymbol{\ell}}_c\|}.
\end{equation}
Then LOS in body and navigation frames is
\begin{equation}
\hat{\mathbf{r}}_b = R_{bc}\hat{\boldsymbol{\ell}}_c,
\qquad
\hat{\mathbf{r}}_n = R_{nb}\hat{\mathbf{r}}_b = R_{nb}R_{bc}\hat{\boldsymbol{\ell}}_c.
\end{equation}
With perfect alignment, $\hat{\mathbf{r}}_b=\hat{\boldsymbol{\ell}}_c$.

Practical note: verify pixel-axis signs and camera convention against FRD (especially image $v$ direction versus body $+Z$ down).

\subsection{PN and ZEM with camera measurements}
\paragraph{PN path (LOS-rate based)}
Estimate LOS angular rate using filtered LOS and gyro compensation. A convenient geometric relation is
\begin{equation}
\boldsymbol{\omega}_{LOS,b} \approx \hat{\mathbf{r}}_b \times \dot{\hat{\mathbf{r}}}_b,
\qquad
\boldsymbol{\omega}_{LOS,n}=R_{nb}\boldsymbol{\omega}_{LOS,b}.
\end{equation}
Then command in navigation frame:
\begin{equation}
\mathbf{a}_{cmd,n}^{PN} = N V_c\left(\boldsymbol{\omega}_{LOS,n}\times\hat{\mathbf{r}}_n\right).
\end{equation}
PN mainly needs reliable LOS/LOS-rate and a practical closing-speed estimate $V_c$.

\paragraph{ZEM path (state-estimate based)}
ZEM requires relative position and velocity estimates:
\begin{equation}
\mathbf{r} = \rho\hat{\mathbf{r}}_n,
\qquad
\mathbf{v}=\mathbf{v}_T-\mathbf{v}_M,
\end{equation}
where range $\rho$ comes from stereo/depth/radar/fused estimator. Then reuse standard ZEM law:
\begin{equation}
t_{go}=\mathrm{clamp}\!\left(\frac{\|\mathbf{r}\|}{\max(V_c,\epsilon)},t_{go,min},t_{go,max}\right),
\end{equation}
\begin{equation}
\mathbf{ZEM}=\mathbf{r}+\mathbf{v}t_{go},
\qquad
\mathbf{a}_{cmd,n}^{ZEM}=N\frac{(I-\hat{\mathbf{r}}\hat{\mathbf{r}}^\top)\mathbf{ZEM}}{t_{go}^2}.
\end{equation}

\subsection{Fixed-Wing realization (bank-to-turn)}
For fixed-wing, use lateral part of $\mathbf{a}_{cmd,n}$ and map to bank command:
\begin{equation}
a_{lat,cmd} = \text{lateral component of }\mathbf{a}_{cmd,n},
\qquad
\phi_{cmd}=\mathrm{sat}\!\left(\arctan\frac{a_{lat,cmd}}{g},-\phi_{max},+\phi_{max}\right).
\end{equation}
Then send roll attitude setpoint (or roll-rate setpoint from an outer roll loop), while longitudinal loops provide pitch/throttle.

\paragraph{Fixed-wing cycle sketch}
\begin{verbatim}
1) Detect target in image -> (u, v)
2) Build LOS unit vector rhat_n from K, R_bc, R_nb
3) Choose guidance:
   - PN: estimate omega_LOS, Vc -> a_cmd_n
   - ZEM: estimate rho, v_rel -> r, v, tgo, ZEM -> a_cmd_n
4) Extract lateral acceleration demand
5) Convert to bank command phi_cmd, apply limits/rate limits
6) Send PX4 roll setpoint; pitch/throttle from longitudinal controllers
\end{verbatim}

\subsection{Multirotor realization (acceleration/attitude control)}
For multirotors, PN/ZEM still generate $\mathbf{a}_{cmd,n}$ in navigation frame; realization differs:
\begin{itemize}
    \item If acceleration setpoints are supported, track $\mathbf{a}_{cmd,n}$ directly with axis saturation.
    \item If attitude/thrust interface is used, convert desired translational acceleration (with gravity compensation) into tilt and thrust commands.
\end{itemize}
Typical form:
\begin{equation}
\mathbf{f}_{des} = m\left(\mathbf{a}_{cmd,n} + g\mathbf{e}_D\right),
\end{equation}
where $\mathbf{e}_D$ is Down-axis unit vector in NED. Desired thrust direction is set by $\mathbf{f}_{des}$, then converted to attitude plus thrust magnitude and clipped by tilt/thrust limits.

Yaw is usually controlled separately (for example, yaw to keep the target near image center) while translational guidance tracks PN/ZEM acceleration demand.

\paragraph{Multirotor cycle sketch}
\begin{verbatim}
1) Detect target in image -> (u, v)
2) Build LOS unit vector rhat_n from K, R_bc, R_nb
3) Choose guidance:
   - PN: estimate omega_LOS, Vc -> a_cmd_n
   - ZEM: estimate rho, v_rel -> r, v, tgo, ZEM -> a_cmd_n
4) Apply accel/tilt/thrust saturation and filtering
5) Convert a_cmd_n (+gravity compensation) to attitude/thrust setpoints
6) Send multirotor setpoints; keep yaw strategy active for camera framing
\end{verbatim}

\subsection*{Estimator requirements summary}
\begin{itemize}
    \item PN with camera can run with LOS and LOS-rate plus approximate $V_c$.
    \item Full ZEM needs range $\rho$ and relative velocity quality; monocular-only setups usually need a fused depth/range estimator.
    \item For both methods, latency compensation and low-pass filtering of LOS/LOS-rate are important for stability.
\end{itemize}

\end{document}
